\setcounter{section}{29}
\section{Zasada nieoznaczoności jako konsekwencja Plancherela}
\label{sec:dt-uncertainty}
\providecommand{\atlasref}[2]{#2}

Por. \atlasref{PDE-08::thm:pde8-uncertainty}{RRC-08: nierówność Heisenberga}.


\begin{twierdzenie}[Nierówność Heisenberga w jednej współrzędnej]
\label{thm:dt-uncertainty}
Jeśli $0\ne f\in\mathcal S(\R^d)$ i $1\le j\le d$, to
\[
 \|x_jf\|_{L^2}\,\|\xi_j\widehat f\|_{L^2}
 \ge\frac12\,\|f\|_{L^2}\,\|\widehat f\|_{L^2}.
\]
Stała $1/2$ jest optymalna. Równość zachodzi
dokładnie wtedy, gdy dla pewnego $\lambda>0$
funkcja $f$ ma postać
\[
 f(x)=e^{-\lambda x_j^2/2}\,g(x'),
\]
gdzie dla $d>1$ przez $x'$ oznaczamy pozostałe
$d-1$ współrzędnych, a $g\in\mathcal S(\R^{d-1})$
jest niezerowa. Dla $d=1$ czynnik $g(x')$
zastępujemy niezerową stałą zespoloną.
\end{twierdzenie}
\begin{proof}
Najpierw zauważmy, że $|f|^2=f\overline f$,
a także $x_j|f|^2$ należą do $\mathcal S$.
Sprzężenie zespolone nie zmienia półnorm.
Dla $u,v\in\mathcal S$ wzór Leibniza daje
\[
 Q_{N,m}(uv)\le 2^m Q_{N,m}(u)Q_{0,m}(v)
 \qquad(N,m\in\N_0),
\]
gdyż każda pochodna iloczynu jest sumą składników
$(\partial^\gamma u)(\partial^{\alpha-\gamma}v)$
ze współczynnikami dwumianowymi o sumie $2^{|\alpha|}\le2^m$.
Iloczyn dwóch funkcji Schwartza należy zatem do
$\mathcal S$, a mnożenie przez współrzędną zachowuje
$\mathcal S$ na mocy stwierdzenia~\ref{prop:schwartz-operators17}.
W szczególności $\partial_j(x_j|f|^2)\in L^1$.
Dla ustalonych pozostałych współrzędnych
podstawowe twierdzenie rachunku całkowego daje
\[
 \int_{-R}^R\partial_j(x_j|f(x)|^2)
 \,\mathrm d\lambda_1(x_j)
 =\left[x_j|f(x)|^2\right]_{x_j=-R}^{x_j=R}.
\]
Prawa strona dąży do zera dla $R\to\infty$,
ponieważ dla $|x_j|=R$ mamy
\[
 |x_j||f(x)|^2
 \le R\,q_{1,0}(f)^2(1+R)^{-2}.
\]
Dla $d>1$ dzięki całkowalności pochodnej
na $\R^d$ twierdzenie Fubiniego zapewnia, że
dla prawie każdego $x'$ pochodna ta jest
bezwzględnie całkowalna na prostej zmiennej $x_j$.
Na takim przekroju całka po $[-R,R]$ zbiega
do całki po $\R$ z twierdzenia o zbieżności
zmajoryzowanej, z majorantą
$|\partial_j(x_j|f|^2)|$ na przekroju.
Z oszacowania składnika brzegowego wynika,
że ta całka po $\R$ jest zerowa. Fubini
pozwala teraz całkować tę równość względem $x'$.
Dla $d=1$ przechodzimy do granicy bezpośrednio,
ponieważ pochodna jest w $L^1(\R)$.
Otrzymujemy
\[
 \int_{\R^d}\partial_j(x_j|f|^2)
 \,\mathrm d\lambda_d(x)=0.
\]
Z reguły Leibniza
\[
 \partial_j(x_j|f|^2)
 =|f|^2+2\operatorname{Re}(x_j\overline f\,\partial_j f).
\]
Oba składniki są całkowalne: dla drugiego
wynika to z nierówności Cauchy'ego--Schwarza
dla $x_jf,\partial_j f\in L^2$.
W konsekwencji
\[
 \begin{aligned}
 \|f\|_2^2
 &=-2\operatorname{Re}
 \int_{\R^d}x_j\overline f(x)\,\partial_j f(x)
 \,\mathrm d\lambda_d(x)\\
 &\le 2\left|
 \int_{\R^d}x_j\overline f(x)\,\partial_j f(x)
 \,\mathrm d\lambda_d(x)\right|\\
 &\le 2\|x_j f\|_2\|\partial_j f\|_2.
 \end{aligned}
\]
Z twierdzeń~\ref{thm:fourier-schwartz17}
i \ref{thm:plancherel22} mamy
\[
 \|\partial_j f\|_2=\|\xi_j\widehat f\|_2,
 \qquad
 \|f\|_2=\|\widehat f\|_2.
\]
Po podstawieniu otrzymujemy
$\|f\|_2^2\le2\|x_jf\|_2\|\xi_j\widehat f\|_2$.
Dzieląc przez $2$ i wykorzystując zachowanie normy
przez unitarną transformatę, otrzymujemy tezę.

Ustalmy teraz warunek równości. Ponieważ
$\|f\|_2>0$, powyższe oszacowanie pokazuje,
że zarówno $x_jf$, jak i $\partial_j f$
mają dodatnie normy $L^2$.
Równość w nierówności Heisenberga jest
równoważna równości w obu nierównościach
w powyższym łańcuchu.
Równość w nierówności Cauchy'ego--Schwarza
oznacza liniową zależność obu wektorów w $L^2$:
$\partial_j f=cx_jf$ dla pewnego $c\in\C$.
Istotnie, dla $p=\partial_j f$, $q=x_jf$
oraz $c=\langle p,q\rangle_2/\|q\|_2^2$
mamy
\[
 \|p-cq\|_2^2
 =\|p\|_2^2-\frac{|\langle p,q\rangle_2|^2}{\|q\|_2^2}.
\]
Równość w nierówności Cauchy'ego--Schwarza
daje zatem $p=cq$ w $L^2$, a zależność ta
jest również wystarczająca dla równości.
Wtedy
\[
 \int_{\R^d}x_j\overline f\,\partial_j f
 \,\mathrm d\lambda_d(x)=c\|x_jf\|_2^2.
\]
Równość w pierwszej nierówności wymaga,
by liczba po lewej stronie była rzeczywista
i ujemna (jej część rzeczywista wynosi
$-\|f\|_2^2/2$). Stąd $c=-\lambda$
dla pewnego $\lambda>0$.
Ponieważ $\partial_j f$ i $x_j f$ są ciągłe,
ich równość prawie wszędzie daje równość
w każdym punkcie: niezerowa różnica w punkcie
byłaby niezerowa na pewnej otwartej kuli
o dodatniej mierze.
Dla ustalonych $x'$ otrzymujemy więc
\[
 \partial_j\bigl(e^{\lambda x_j^2/2}f(x)\bigr)=0,
\]
czyli $f(x)=e^{-\lambda x_j^2/2}g(x')$,
gdzie $g(x')=f(x)|_{x_j=0}$.
Dla $d>1$ funkcja $g$ jest gładka jako
ograniczenie funkcji gładkiej do hiperpłaszczyzny.
Jeśli $\alpha',\beta'\in\N_0^{d-1}$,
to przez $\alpha,\beta\in\N_0^d$ oznaczmy
wielowskaźniki powstałe przez wstawienie
zera na $j$-tej pozycji.
Wtedy
\[
 \sup_{x'\in\R^{d-1}}
 |(x')^{\alpha'}\partial^{\beta'}g(x')|
 \le p_{\alpha,\beta}(f)<\infty.
\]
Zatem $g\in\mathcal S(\R^{d-1})$ i $g\ne0$,
bo w przeciwnym razie wzór na $f$ dawałby
$f=0$. Dla $d=1$ tym samym argumentem
otrzymujemy niezerową stałą $f(0)$.

Odwrotnie, dla dowolnej $\lambda>0$ i takiej
niezerowej $g$ funkcja
$f(x)=e^{-\lambda x_j^2/2}g(x')$ należy
do $\mathcal S$. Dla $d>1$, oznaczając
$G(t)=e^{-\lambda t^2/2}$, mamy
\[
 p_{\alpha,\beta}(f)=
 p_{\alpha_j,\beta_j}(G)\,
 p_{\alpha',\beta'}(g)<\infty,
\]
gdzie $\alpha',\beta'$ powstają przez
usunięcie $j$-tej współrzędnej.
Wynika to z rozdzielenia pochodnych i wag
oraz supremum niezależnych nieujemnych czynników.
Dla $d=1$ wystarcza przynależność Gaussa do $\mathcal S$.
Niezerowość $f$ wynika z niezerowości $g$.
Spełnia $\partial_j f=-\lambda x_j f$,
więc obie nierówności w dowodzie stają się
równościami. Takie funkcje istnieją w każdym
wymiarze, na przykład $f(x)=e^{-|x|^2/2}$.
Zatem stałej $1/2$ nie można zwiększyć.
\end{proof}

\begin{wniosek}[Nierówność z przesuniętymi środkami]
Dla $f\in\mathcal S(\R^d)$ i dowolnych $a,b\in\R$
mamy
\[
 \|(x_j-a)f\|_2\,\|(\xi_j-b)\widehat f\|_2
 \ge\tfrac12\|f\|_2\|\widehat f\|_2.
\]
\end{wniosek}
\begin{proof}
Dla $e_j$ oznaczającego $j$-ty wektor
bazy standardowej połóżmy
$h(x)=e^{-ibx_j}f(x+ae_j)$.
Przesunięcie zachowuje $\mathcal S$ przez
nierówność $1+|x|\le(1+|a|)(1+|x+ae_j|)$,
a mnożenie przez $e^{-ibx_j}$ zachowuje
$\mathcal S$, gdyż wszystkie pochodne
mnożnika są ograniczone.
Ponieważ $h\in\mathcal S\subset L^1$,
podstawienie $y=x+ae_j$ w całce Fouriera daje
\[
 \begin{aligned}
 \widehat h(\xi)
 &=(2\pi)^{-d/2}\int_{\R^d}e^{-ix\cdot(\xi+be_j)}
 f(x+ae_j)\,\mathrm d\lambda_d(x)\\
 &=e^{ia(\xi_j+b)}(2\pi)^{-d/2}
 \int_{\R^d}e^{-iy\cdot(\xi+be_j)}f(y)
 \,\mathrm d\lambda_d(y)\\
 &=e^{ia(\xi_j+b)}\widehat f(\xi+be_j).
 \end{aligned}
\]
Czynniki wykładnicze mają moduł $1$,
a miara Lebesgue'a jest niezmiennicza
względem przesunięć. Stąd
\[
 \begin{aligned}
 \|h\|_2&=\|f\|_2,&
 \|x_jh\|_2&=\|(x_j-a)f\|_2,\\
 \|\widehat h\|_2&=\|\widehat f\|_2,&
 \|\xi_j\widehat h\|_2&=\|(\xi_j-b)\widehat f\|_2.
 \end{aligned}
\]
Dla $f\ne0$ stosujemy twierdzenie
~\ref{thm:dt-uncertainty} do $h$;
dla $f=0$ nierówność jest oczywista.
\end{proof}

Dla $f\ne0$ można w szczególności wybrać
\[
 a=\frac{\int_{\R^d}x_j|f(x)|^2\,
 \mathrm d\lambda_d(x)}{\|f\|_2^2},
 \qquad
 b=\frac{\int_{\R^d}\xi_j|\widehat f(\xi)|^2\,
 \mathrm d\lambda_d(\xi)}{\|\widehat f\|_2^2}.
\]
Całki są bezwzględnie zbieżne, ponieważ
\[
 \begin{aligned}
 \int_{\R^d}|x_j||f(x)|^2\,\mathrm d\lambda_d(x)
 &\le\|x_jf\|_2\|f\|_2,\\
 \int_{\R^d}|\xi_j||\widehat f(\xi)|^2\,
 \mathrm d\lambda_d(\xi)
 &\le\|\xi_j\widehat f\|_2\|\widehat f\|_2.
 \end{aligned}
\]
Oba mianowniki są dodatnie z niezerowości $f$
i twierdzenia Plancherela. Liczby $a,b$ są
rzeczywiste, a obie niżej wymienione gęstości
mają całkę $1$ z definicji normy $L^2$.
Drugie momenty są skończone, ponieważ
$x_jf,\xi_j\widehat f\in L^2$.
Jeśli $\sigma_x^2$ i $\sigma_\xi^2$
oznaczają wariancje $j$-tej współrzędnej
względem miar probabilistycznych o gęstościach
$|f|^2/\|f\|_2^2$ oraz
$|\widehat f|^2/\|\widehat f\|_2^2$, to
\[
 \sigma_x=\frac{\|(x_j-a)f\|_2}{\|f\|_2},
 \qquad
 \sigma_\xi=\frac{\|(\xi_j-b)\widehat f\|_2}
 {\|\widehat f\|_2}.
\]
Z poprzedniego wniosku dostajemy
$\sigma_x\sigma_\xi\ge1/2$.
Stała odpowiada konwencji z definicji~\ref{def:fourier-schwartz17},
w której wykładnik transformaty to $-ix\cdot\xi$.
